\[
M_m:=F_{m+2},
\qquad
Z_m:=\left\{k\in \mathbb Z/(M_m\mathbb Z):\widehat d_m(k)=0\right\},
\qquad
\Gamma_m:=M_m\left(\operatorname{Col}_m-\frac1{M_m}\right).
\]

\begin{theorem}[零点完全空缺而缩放碰撞缺口仍爆发]\label{thm:conclusion-fold-zerofree-subsequence-gap-explosion}
设
\[
\mathcal N:=\{m\ge 1:3\nmid (m+2)\}.
\]
则对每个 \(m\in \mathcal N\) 都有
\[
Z_m=\varnothing,
\]
但沿同一子列仍成立
\[
\Gamma_m\longrightarrow +\infty
\qquad
(m\to\infty,\ m\in\mathcal N).
\]
因此，谱零点完全消失与碰撞缺口可忽略在该框架内是严格不相干的两件事。
\end{theorem}
\begin{proof}
主稿关于谱零点的判别已经给出：若
\[
3\nmid (m+2),
\]
则
\[
Z_m=\varnothing.
\]
该条件定义的子列 \(\mathcal N\) 显然无限。

另一方面，定理 \ref{thm:conclusion-pisot-bernoulli-nonrajchman-collision-inflation} 已证明
\[
F_{m+2}\operatorname{Col}_m\longrightarrow +\infty.
\]
于是
\[
\Gamma_m
=
F_{m+2}\operatorname{Col}_m-1
\longrightarrow
+\infty
\]
沿全序列成立，故沿子列 \(\mathcal N\) 亦成立。证毕。
\end{proof}
\leanverified{paper\_conclusion\_fold\_zerofree\_subsequence\_gap\_explosion}

\begin{corollary}[零点计数不可能渐近控制缩放碰撞缺口]\label{cor:conclusion-fold-zero-count-cannot-control-gap}
不存在函数
\[
\omega:[0,1]\to[0,\infty)
\]
与数列 \(\varepsilon_m\downarrow 0\)，使得
\[
\Gamma_m\le \omega\!\left(\frac{|Z_m|}{F_{m+2}}\right)+\varepsilon_m
\]
对所有充分大的 \(m\) 成立，并同时满足
\[
\omega(t)\to 0
\qquad
(t\downarrow 0).
\]
\end{corollary}
\begin{proof}
若存在这样的 \(\omega\) 与 \(\varepsilon_m\)，则沿定理 \ref{thm:conclusion-fold-zerofree-subsequence-gap-explosion} 的子列 \(\mathcal N\) 有
\[
\omega\!\left(\frac{|Z_m|}{F_{m+2}}\right)+\varepsilon_m
=
\omega(0)+\varepsilon_m
\longrightarrow 0,
\]
因为此时 \(|Z_m|=0\) 且 \(\omega(t)\to 0\) 当 \(t\downarrow 0\)。但同一定理又给出
\[
\Gamma_m\to+\infty
\qquad
(m\to\infty,\ m\in\mathcal N),
\]
矛盾。证毕。
\end{proof}
